% CS 418: Intro to Data Science — Worksheet 12 (12-week06-day2)
% Week 6, Wednesday 2026-09-30
% NOTICE: This worksheet was drafted by an LLM coding system from the
% instructor's list of checkpoints. It is maintained, reviewed, and vetted by humans.
% Compile with: latexmk -pdf worksheet.tex

\documentclass[11pt]{article}

% ── Packages ──────────────────────────────────────────────────────────────────
\usepackage[margin=0.65in, headheight=14pt]{geometry}
\usepackage{amsmath, amssymb}
\usepackage{ragged2e}
% Never break a word across lines, including at hyphens ("Sans-Serif").
\hyphenpenalty=10000
\exhyphenpenalty=10000
\usepackage{enumitem}
\usepackage{booktabs}
\usepackage{graphicx}
\usepackage{hyperref}
\usepackage{xcolor}
\usepackage{fancyhdr}
\usepackage{mdframed}
\usepackage{array}

% ── Page style ────────────────────────────────────────────────────────────────
\pagestyle{fancy}
\fancyhf{}
\lhead{\textbf{CS 418: Intro to Data Science}}
\rhead{UIC, Fall 2026}
\cfoot{\thepage}
\renewcommand{\headrulewidth}{1pt}
\setlength{\headheight}{12pt}
\setlength{\headsep}{0.4cm}

% ── Hyperlinks ────────────────────────────────────────────────────────────────
\hypersetup{colorlinks=false, hidelinks}

% ── Convenience environments ──────────────────────────────────────────────────
% Usage: \blankspace{6cm} — blank area for hand-written work, no rules
\newcommand{\blankspace}[1]{%
  \vspace{0.3em}
  \begin{minipage}[t][#1][t]{\linewidth}
  \end{minipage}
  \vspace{0.3em}
}

% Usage: \guidedopen{title}{guided content}{guided workspace height}{open-ended workspace height}
% Left: a titled, structured prompt with its own blank workspace. Right: a
% fully blank workspace for open-ended exploration of the same topic (no
% prompt text), sized so the two columns land at roughly the same total
% height. Neither workspace is ruled.
\newcommand{\guidedopen}[4]{%
  \noindent\textbf{#1}\\[0.18em]
  \noindent
  \begin{minipage}[t]{0.485\textwidth}
    #2
    \blankspace{#3}
  \end{minipage}
  \hfill
  \begin{minipage}[t]{0.485\textwidth}
    \blankspace{#4}
  \end{minipage}
  \\[0.4em]
}

% Usage: \panelbox{width}{height}{caption} — a framed drawing panel with `caption`
% centered beneath it. An empty caption leaves the panel bare, with nothing
% below the frame at all.
\newcommand{\panelbox}[3]{%
  \begin{minipage}[t]{\dimexpr#1+2\fboxsep+2\fboxrule\relax}
    \framebox[#1][c]{\rule{0pt}{#2}}%
    \if\relax\detokenize{#3}\relax\else
      \\[0.25em]{\centering\small #3\par}%
    \fi
  \end{minipage}%
}

% ══════════════════════════════════════════════════════════════════════════════
\begin{document}
\RaggedRight

% ── Title block: topic left, info box right ───────────────────────────────────
\noindent
\begin{minipage}[t]{0.55\textwidth}
  \vspace{0pt}
  {\LARGE \textbf{Worksheet \#\,12}} \\[0.18em]
  {\large \textit{Week 6, Day 2: Visualization rhetoric and causality}}
\end{minipage}
\hfill
\begin{minipage}[t]{0.40\textwidth}
  \vspace{0pt}
  \small
  Name: \underline{\hspace{4.5cm}} \\[0.35em]
  NetID: \underline{\hspace{4.3cm}} \\[0.35em]
  Date: \underline{\hspace{4.4cm}} \\
\end{minipage}

\vspace{0.6em}
\noindent\rule{\linewidth}{0.6pt}
\vspace{0.1em}

\noindent\underline{\textbf{Guided checkpoints}}\\[0.2em]
\noindent{\small For full credit, fill in the checkpoints. If you skip a checkpoint, you must include your own notes for full credit.}
\vspace{0.8em}

% ── 1. Group checkpoint ──────────────────────────────────────────────────────
\guidedopen{1. Group Checkpoint}{
(a) Does your group have a name yet? If so, what is it? If not, list a few ideas.

\vspace{2.4cm}

(b) What will \textit{you} do next for your group project?
}{2.6cm}{7.2cm}

\vfill

% ── 2. Misleading axes ───────────────────────────────────────────────────────
\guidedopen{2. Misleading Axes}{
List three axis choices that can be misleading.
}{3.4cm}{4.3cm}

\vfill

% ── 3. Covariates for ridership ──────────────────────────────────────────────
\guidedopen{3. Covariates for Ridership}{
Discuss with a neighbor: what ``covariates'' could help predict daily ridership at UIC-Halsted?
}{3.6cm}{5.5cm}

\vfill

\newpage

% ── 4. DAG for ridership ─────────────────────────────────────────────────────
\guidedopen{4. A DAG for Ridership}{
Draw a DAG for CTA ridership.
}{4.2cm}{5.1cm}

\vfill

% ── 5. Elemental confounds ───────────────────────────────────────────────────
\guidedopen{5. The Elemental Confounds}{
Draw each confound with $X$, $Y$, and $Z$.

\vspace{0.45em}

\noindent
\panelbox{7.8cm}{2.8cm}{fork}

\vspace{0.5em}

\noindent
\panelbox{7.8cm}{2.8cm}{pipe}

\vspace{0.5em}

\noindent
\panelbox{7.8cm}{2.8cm}{collider}

\vspace{0.5em}

\noindent
\panelbox{7.8cm}{2.8cm}{descendant}
}{0.2cm}{14.5cm}

\vfill

\end{document}
